\begin{beginnerstatement}[Kurz erklärt: reproduzierbare Ergebnisse]

Ein reproduzierbarer Build soll bei gleichen Eingaben wieder dasselbe Ergebnis
erzeugen. Für diese Fallstudie wurde jedoch nur gezeigt, dass zwei
Generierungen von \texttt{desktop-cloud + postgres} auf
\texttt{1cdfb6e} in derselben Umgebung bytegleich waren. Eine festgelegte
Dependency-Version bestimmt genau, welcher Stand einer benötigten Bibliothek
verwendet wird. Ein Lockfile schreibt solche Versionen fest, statt beim
nächsten Lauf beliebige neuere Pakete zu wählen. Ein Hash ist ein kurzer
digitaler Fingerabdruck, mit dem sich ein bestimmter Inhalt sehr zuverlässig
prüfen lässt. Die Toolchain umfasst die eingesetzten Laufzeiten, Compiler und
Buildwerkzeuge. Unterschiede bei Betriebssystem, Toolchain, Netzwerk oder
Plattform-SDKs können die Wiederholung trotzdem beeinflussen.

\end{beginnerstatement}
